\scanpage{114}
\uline{Supposons} maintenant que l'on ait $\sum\alpha_i<+\infty$,
posons $f(z)=\prod_1^\infty(1+\alpha_i z)$. Alors $\rho$ est aussi
l'ordre de la fonction entière, i.e
\[ \rho=\limsup_{0<r\to\infty}\frac{\log(\log f(r))}{\log r} \]

En effet, on sait que $\rho$ est $\leq$ ordre de la
fonction $f(z)$, \deleted{\unsure{et}} \unsure{montrons|démontrons} le contraire. On est
\uline{ramené} à montrer que \unsure{.|:} \deleted{le produit} dans le
cas où on a $\alpha_n=\frac1{n^\beta}$, on a
$\displaystyle \limsup_{0<r\to\infty}\frac{\log\log f(r)}{\log r}\leq\frac1\beta$ \quad (en somme on a en fait $=\frac1\beta$)

En effet, on a
\[ f(r)=\prod_1^\nu\left(1+\frac1{n^\beta}r\right)\prod_{\nu+1}^{\infty}\left(1+\frac1{n^\beta}r\right) \]
le deuxième facteur est $\leq e^{R(\nu)r}$
avec $\displaystyle R(\nu)=\sum_{\nu+1}^\infty\frac1{n^\beta}\leq\int_\nu^\infty\frac{du}{x^\beta}=\frac1{\beta-1}\frac1{\nu^{\beta-1}}$
Prenons $\nu$ tel que $r\unsure{\cdot}\frac1{\nu^{\beta-1}}\leq r^{\frac1\beta}$, i.e $\nu\geq r^{\frac1\beta}$
de façon précise, $\nu=[r^{\frac1\beta}]+1$. Alors, $e^{R(\nu)r}$
sera une \uline{fonction de $r$} d'ordre $\leq\frac1\beta$\unsure{;|.} Il
suffit de le vérifier pour le produit $\prod_1^\nu(1+\frac1{n^\beta}r)$
$\leq(1+r)^\nu=e^{\nu\log(1+r)}$, qui est manifestement
d'ordre $\leq\frac1\beta$.

Rappelons que si on pose $f(\unsure{r|z})=\sum a_n z^n$,
\deleted{\unsure{dans|donc}} et $\unsure{B|D}(r)=\mathop{\rm Sup}_n(a_n r^n)$, alors on
a aussi \qquad ordre \unsure{a|de} $f=\limsup_{r\to\infty}\frac{\log\log B(r)}{\log r}$
\margin{115}
